\documentclass{zhvt-classic}

% 使用本地源流明体（GenRyuMin2 TW）按文件路径加载，不依赖系统字体名
% 版框（内框／界栏／书口侧栏／鱼尾）改用与版心印章同一印泥棕色：
%   印章图片 yshs.png 实色 = RGB(135,68,52) = HTML 874434
\usepackage{xcolor}
\usepackage{tikz}
% sealbrown 双注册：xcolor 供 TikZ/pgf，L3 供 cls 版框
\definecolor{sealbrown}{HTML}{874434}
\ExplSyntaxOn
\color_set:nnn { sealbrown } { HTML } { 874434 }
\ExplSyntaxOff

\zhvtset{
  grid color       = sealbrown,
  main font        = GenRyuMin2TW-SB,
  main font path   = C:/Users/Administrator/texmf/fonts/opentype/public/GenRyuMin2TW/,
jiazhu font      = GenRyuMin2TW-M,
  jiazhu font path = C:/Users/Administrator/texmf/fonts/opentype/public/GenRyuMin2TW/,
  % 页码框容纳 4 字：正常页「九十七」+前置空 1 格=4 行恰好；
  % 100+ 页「一百零一」共 5 行超出，末字「一/二/三」被框裁掉（2026-10-04 实测），故放宽到 5
  shukou page number chars = 5,
  % 书口书名框高从默认 6 字减到 4 字，把省出的 2 字高拨给 chapter 框，
  % 使 6 字书名+卷名放进 9 字高的 chapter 框后不再压线（2026-10-05）
  shukou title chars = 4,
}

\title{史记}[郑世家]
\author{汉·司马迁}
\maker{闲叟漫记·重製}


% ===== 索隐述赞样式（2026-11-08 四次定稿 · 愚大人亲手改定参数）=====
% ⚠ 竖排前提：自定义字体家族必须带 RawFeature={vertical:+vert}；
%   simhei.ttf 加 +vhal 会整体加载失败静默回退，只写 +vert；simfang 同理。
\newCJKfontfamily{\shuzanfont}[Path=C:/Windows/Fonts/,RawFeature={vertical:+vert}]{simhei.ttf}
\newCJKfontfamily{\shuzanbody}[Path=C:/Windows/Fonts/,RawFeature={vertical:+vert}]{simfang.ttf}
% 题头：152x30pt 圆角 6pt 棕底白字；baseline (0,2) 使整框落入界行格子不压版框线；
% 白字节点 x 不补偿、y +8.2pt（像素级墨迹实测 dx/dy≈0，L/R 留白约 20pt）。
\newcommand{\shuzanhead}[1]{%
  \begin{tikzpicture}[baseline={(0,2)}]
    \fill[sealbrown, rounded corners=6pt]
      (-76pt,-15pt) rectangle (76pt,15pt);
    \node[text=white, inner sep=0pt, anchor=center] at (0pt,8.2pt)
      {\shuzanfont\fontsize{23}{23}\selectfont #1};
  \end{tikzpicture}%
}
% ★ 修正书口卷名位置：chapter 框起始位置改回 cls 默认值，避免全局影响。
%   只在每个卷首页通过 \juan 命令里的空命令微调该页位置（安全，不会跨页污染）。
\ExplSyntaxOn
\AtBeginDocument{
  \dim_set:Nn \l__zhvt_shukou_chapter_title_start_dim
    { \l__zhvt_shukou_title_height_dim + \l__zhvt_font_size_dim * 2 }
  \dim_set:Nn \l__zhvt_shukou_page_number_start_dim
    { \l__zhvt_shukou_chapter_title_start_dim + \l__zhvt_shukou_chapter_title_height_dim }
}
\ExplSyntaxOff

\begin{document}

% ★★ 版心（书口）排布覆写：把书名放到鱼尾之下（2026-10-04）
%
%   cls 默认版心自上而下：书名（title 静态框，在鱼尾之上）→ 鱼尾 → 篇名（chapter 动态框）→ 页码
%   本模板改为：          鱼尾（位置不动）→ 书名 → 篇名 → 页码分割线 → 页码
%
%   手法：把 cls 放在 title 静态框里的书名，改放进 chapter 动态框（鱼尾之下）；
%         title 静态框留空。书名取 \l__zhvt_title_str，即 \title{...} 的第一参数，
%         所以换书时只改 \title，不必在此重复写书名。
%
%   ⚠ 三条必须遵守的写法（2026-10-04 实测）：
%     ① 汉字必须用 \hbox{...} 包住。flowfram 会把框内容装进一个宽 14.4mm（= 一个 line_sep
%        = 书口栏宽）的纵向 minipage 并按该宽度折行，而整页靠 PDF /Rotate 90 显示、
%        框内文字本身不变向 —— 不包 \hbox 时「围炉夜话」4 字会被逐字折成 4 行，
%        在版面上就成了横排 4 字并横向溢出版框。
%     ② 竖排下**一个 \hbox 就是一整列**。书名与篇名要同处一列，必须放进同一个 \hbox；
%        分成两个 \hbox 时，后一个会被当成新的一列、甩到隔壁栏去
%        （实测「序」跑到书口栏外侧 14.4mm 的正文列上）。中间的 \quad 在框内是横向空白，
%        显示出来正好是书名与篇名之间的纵向间隔。
%     ③ ★ 内容盒要写成 **与框等宽的定宽盒** `\hbox to \l__zhvt_shukou_chapter_title_height_dim{…\hfil}`。
%        实测：chapter 框是把内容盒沿纵向**居中**摆放的，内容宽度一变，书名就跟着漂：
%        自然宽度的 \hbox（书名 38.4mm）落在 x=157.2mm，加篇名后（57.6mm）落在 x=152.4mm，
%        目录页（无篇名）与正文页的书名就差半格。做成等宽盒后，无论有没有篇名，
%        内容固定从框顶起排，实测两页均落在 x=142.8mm，不再漂。
\ExplSyntaxOn
% 把书口 chapter 框（书名+篇名）起始位置下移 1 字高，避免 4 字卷名顶到鱼尾
% ⚠ cls 的 \__zhvt_calculate_dimensions: 在 \AtBeginDocument 里执行，会重算该 dim；
%   必须把覆写挂到 begin document 之后，否则会被二次覆盖（2026-10-05 实测）。
% (旧方案已移到导言区，直接修改 chapter 框起始位置)
% 卷标题命令 #1=书口篇名(鱼尾) #2=正文卷名（含卷序，卷名内空格用 ~）
%
% ★★ 目录页两项修订（2026-10-05，愚大人指出）★★
%   ① 页码起点统一对齐：
%      原状 —— 目录项写成 \par\hspace{4em}卷名\hspace{4em}页码，页码起点 = 卷名实宽 + 4em。
%      卷名实宽随字数变：「卷一~梁惠王上」6 字 +~、「卷七~离娄上」5 字 +~，
%      差 1 个字位 → 卷七~卷十 4 条的页码整体比其余 10 条高 27.21pt（一个字），看着上下不齐。
%      改法 —— 卷名补全角空格（宽 = 1 字 = 27.215pt）补足到 **7 字符位**再写进目录，
%      14 条页码起点遂全部落在同一条竖线上（起点对齐；末字随页码字数自然参差）。
%      ⚠ 只对目录里的副本 \l_tmpc_tl 补，正文卷名仍用 \l_tmpa_tl 原样 —— 正文像素不受影响。
%   ② 卷名整体上移半字位：前导 \hspace 由 4em 改 3.5em。
%      实测原状目录块在版框内偏下（天头留白 139.6pt / 地脚留白 110.7pt），
%      上移 0.5em = 13.6pt 后上下基本均分，视觉居中。
\NewDocumentCommand { \juan } { mm } {
	\clearpage
	\tl_set:Nn \l__zhvt_chapter_title_tl { #1 }
	\tl_set:Nn \l_tmpa_tl { #2 }
	\tl_set_eq:NN \l_tmpc_tl \l_tmpa_tl
	\int_step_inline:nn { 7 - \str_count:n { #2 } }
	  { \tl_put_right:Nn \l_tmpc_tl { 　 } }
	\tl_set:Nx \l_tmpb_tl { \zhnum { page } }
	\tl_put_left:Nn \l_tmpb_tl { \hspace { 4em } }
	\tl_put_left:Nn \l_tmpb_tl { \l_tmpc_tl }
	\tl_put_left:Nn \l_tmpb_tl { \par \hspace { 3.5em } }
	\addtocontents { toc } { \l_tmpb_tl }
	\skip_horizontal:n { \l__zhvt_font_size_dim * 4 }
	\l_tmpa_tl
	\par
	\thispagestyle { empty }
}
\cs_set_protected:Npn \__zhvt_set_shukou_content: {
	% ★ 鱼尾上方留白（title 静态框）嵌入闲叟钤印：
	%   整页 /Rotate 90，故此处把图用 \rotatebox{90} 转正为竖章；
	%   尺寸卡在版心条幅内（书内约 7mm 宽 × 33.8mm 高，留余量不压鱼尾）。
	\setstaticcontents * { title } {
		\vbox_to_ht:nn { \l__zhvt_baseline_skip_dim } {
			\vskip 14.2pt
			\hbox_to_wd:nn { \l__zhvt_shukou_title_height_dim } {
				\hfil
				\rotatebox [ origin = c ] { 90 } { \includegraphics [ width = 7mm , height = 33.8mm ] { yshs.png } }
				\hfil
			}
		}
	}
% ★ 版心（书口）统一改法（2026-11-16 愚大人定）：18pt＋框内居中＋书口文字同版框色
%   ① 原 27.2pt 时 7 字长篇名（如「廉颇蔺相如列传」）末字溢出 chapter 框、压住页码；
%      改 \fontsize{18pt}{18pt}，前后 \hfil 令字串在框内居中（不随篇名字数漂移）。
%   ② \vbox{\vskip 8.2pt} 定量下移，抵消缩字号后栏宽方向 8.2pt 上偏（像素实测 y 中心 543.3 = 原版）。
%      ⚠ \raisebox / \vspace* 在 flowfram 框内无效；换字号须同步换补偿：16pt→10.1、19pt→7.4。
%   ③ chapter 与 page~number 均置 \color{sealbrown}（HTML 874434），与版框／界行同色；正文仍黑。
%   ⚠ 页码字位、钤印、鱼尾、界行、正文一律未动；五部合订本同此口径。
	\setdynamiccontents * { chapter }
	{ \color{sealbrown} \vbox { \vskip 8.2pt \hbox to \l__zhvt_shukou_chapter_title_height_dim
		{ \hfil \fontsize{18pt}{18pt}\selectfont \l__zhvt_title_str \quad \l__zhvt_chapter_title_tl \hfil } } }
	\setdynamiccontents * { page~number } { \color{sealbrown} \quad \zhnum { page } }
}
\__zhvt_set_shukou_content:
\ExplSyntaxOff

\maketitle

% 单篇，无多卷目录，不渲染目录页

\juan{郑世家}{卷四十二~郑世家第十二}
\hspace*{2em}郑桓公友者。周厉王少子而宣王庶弟也。宣王立二十二年。友初封于郑。封三十三岁。百姓皆便爱之。幽王以为司徒。和集周民。周民皆说。河雒之间。人便思之。为司徒一岁。幽王以襃后故。王室治多邪。诸侯或畔之。於是桓公问太史伯曰。﹁王室多故。予安逃死乎。﹂太史伯对曰。﹁独雒之东土。河济之南可居。﹂公曰。﹁何以。﹂对曰。﹁地近虢。郐。虢。郐之君贪而好利。百姓不附。今公为司徒。民皆爱公。公诚请居之。虢。郐之君见公方用事。轻分公地。公诚居之。虢。郐之民皆公之民也。﹂公曰。﹁吾欲南之江上。何如。﹂对曰。﹁昔祝融为高辛氏火正。其功大矣。而其於周未有兴者。楚其後也。周衰。楚必兴。兴。非郑之利也。﹂公曰。﹁吾欲居西方。何如。﹂对曰。﹁其民贪而好利。难久居。﹂公曰。﹁周衰。何国兴者。﹂对曰。﹁齐。秦。晋。楚乎。夫齐。姜姓。伯夷之後也。伯夷佐尧典礼。秦。嬴姓。伯翳之後也。伯翳佐舜怀柔百物。及楚之先。皆尝有功於天下。而周武王克纣後。成王封叔虞于唐。其地阻险。以此有德与周衰并。亦必兴矣。﹂桓公曰。﹁善。﹂於是卒言王。东徙其民雒东。而虢。郐果献十邑。竟国之。\\
\hspace*{2em}二岁。犬戎杀幽王於骊山下。并杀桓公。郑人共立其子掘突。是为武公。\\
\hspace*{2em}武公十年。娶申侯女为夫人。曰武姜。生太子寤生。生之难。及生。夫人弗爱。後生少子叔段。段生易。夫人爱之。二十七年。武公疾。夫人请公。欲立段为太子。公弗听。是岁。武公卒。寤生立。是为庄公。\\
\hspace*{2em}庄公元年。封弟段於京。号太叔。祭仲曰。﹁京大於国。非所以封庶也。﹂庄公曰。﹁武姜欲之。我弗敢夺也。﹂段至京。缮治甲兵。与其母武姜谋袭郑。二十二年。段果袭郑。武姜为内应。庄公发兵伐段。段走。伐京。京人畔段。段出走鄢。鄢溃。段出奔共。於是庄公迁其母武姜於城颍。誓言曰。﹁不至黄泉。毋相见也。﹂居岁馀。已悔思母。颍谷之考叔有献於公。公赐食。考叔曰。﹁臣有母。请君食赐臣母。﹂庄公曰。﹁我甚思母。恶负盟。柰何。﹂考叔曰。﹁穿地至黄泉。则相见矣。﹂於是遂从之。见母。\\
\hspace*{2em}二十四年。宋缪公卒。公子冯奔郑。郑侵周地。取禾。二十五年。卫州吁弑其君桓公自立。与宋伐郑。以冯故也。二十七年。始朝周桓王。桓王怒其取禾。弗礼也。二十九年。庄公怒周弗礼。与鲁易祊。许田。三十三年。宋杀孔父。三十七年。庄公不朝周。周桓王率陈。蔡。虢。卫伐郑。庄公与祭仲。高渠弥发兵自救。王师大败。祝聸射中王臂。祝聸请从之。郑伯止之。曰。﹁犯长且难之。况敢陵天子乎。﹂乃止。夜令祭仲问王疾。\\
\hspace*{2em}三十八年。北戎伐齐。齐使求救。郑遣太子忽将兵救齐。齐釐公欲妻之。忽谢曰。﹁我小国。非齐敌也。﹂时祭仲与俱。劝使取之。曰。﹁君多内宠。太子无大援将不立。三公子皆君也。﹂所谓三公子者。太子忽。其弟突。次弟子亹也。\\
\hspace*{2em}四十三年。郑庄公卒。初。祭仲甚有宠於庄公。庄公使为卿。公使娶邓女。生太子忽。故祭仲立之。是为昭公。\\
\hspace*{2em}庄公又娶宋雍氏女。生厉公突。雍氏有宠於宋。宋庄公闻祭仲之立忽。乃使人诱召祭仲而执之。曰。﹁不立突。将死。﹂亦执突以求赂焉。祭仲许宋。与宋盟。以突归。立之。昭公忽闻祭仲以宋要立其弟突。九月丁亥。忽出奔卫。己亥。突至郑。立。是为厉公。\\
\hspace*{2em}厉公四年。祭仲专国政。厉公患之。阴使其婿雍纠欲杀祭仲。纠妻。祭仲女也。知之。谓其母曰。﹁父与夫孰亲。﹂母曰。﹁父一而已。人尽夫也。﹂女乃告祭仲。祭仲反杀雍纠。戮之於市。厉公无柰祭仲何。怒纠曰。﹁谋及妇人。死固宜哉。﹂夏。厉公出居边邑栎。祭仲迎昭公忽。六月乙亥。复入郑。即位。\\
\hspace*{2em}秋。郑厉公突因栎人杀其大夫单伯。遂居之。诸侯闻厉公出奔。伐郑。弗克而去。宋颇予厉公兵。自守於栎。郑以故亦不伐栎。\\
\hspace*{2em}昭公二年。自昭公为太子时。父庄公欲以高渠弥为卿。太子忽恶之。庄公弗听。卒用渠弥为卿。及昭公即位。惧其杀己。冬十月辛卯。渠弥与昭公出猎。射杀昭公於野。祭仲与渠弥不敢入厉公。乃更立昭公弟子亹为君。是为子亹也。无谥号。\\
\hspace*{2em}子亹元年七月。齐襄公会诸侯於首止。郑子亹往会。高渠弥相。从。祭仲称疾不行。所以然者。子亹自齐襄公为公子之时。尝会斗。相仇。及会诸侯。祭仲请子亹无行。子亹曰。﹁齐彊。而厉公居栎。即不往。是率诸侯伐我。内厉公。我不如往。往何遽必辱。且又何至是。﹂卒行。於是祭仲恐齐并杀之。故称疾。子亹至。不谢齐侯。齐侯怒。遂伏甲而杀子亹。高渠弥亡归。归与祭仲谋。召子亹弟公子婴於陈而立之。是为郑子。是岁。齐襄公使彭生醉拉杀鲁桓公。\\
\hspace*{2em}郑子八年。齐人管至父等作乱。弑其君襄公。十二年。宋人长万弑其君湣公。郑祭仲死。\\
\hspace*{2em}十四年。故郑亡厉公突在栎者使人诱劫郑大夫甫假。要以求入。假曰。﹁舍我。我为君杀郑子而入君。﹂厉公与盟。乃舍之。六月甲子。假杀郑子及其二子而迎厉公突。突自栎复入即位。初。内蛇与外蛇斗於郑南门中。内蛇死。居六年。厉公果复入。入而让其伯父原曰。﹁我亡国外居。伯父无意入我。亦甚矣。﹂原曰。﹁事君无二心。人臣之职也。原知罪矣。﹂遂自杀。厉公於是谓甫假曰。﹁子之事君有二心矣。﹂遂诛之。假曰。﹁重德不报。诚然哉。﹂\\
\hspace*{2em}厉公突後元年。齐桓公始霸。\\
\hspace*{2em}五年。燕。卫与周惠王弟穨伐王。王出奔温。立弟穨为王。六年。惠王告急郑。厉公发兵击周王子穨。弗胜。於是与周惠王归。王居于栎。七年春。郑厉公与虢叔袭杀王子穨而入惠王于周。\\
\hspace*{2em}秋。厉公卒。子文公踕立。厉公初立四岁。亡居栎。居栎十七岁。复入。立七岁。与亡凡二十八年。\\
\hspace*{2em}文公十七年。齐桓公以兵破蔡。遂伐楚。至召陵。\\
\hspace*{2em}二十四年。文公之贱妾曰燕姞。梦天与之兰。曰。﹁余为伯鯈。余。尔祖也。以是为而子。兰有国香。﹂以梦告文公。文公幸之。而予之草兰为符。遂生子。名曰兰。\\
\hspace*{2em}三十六年。晋公子重耳过。文公弗礼。文公弟叔詹曰。﹁重耳贤。且又同姓。穷而过君。不可无礼。﹂文公曰。﹁诸侯亡公子过者多矣。安能尽礼之。﹂詹曰。﹁君如弗礼。遂杀之。弗杀。使即反国。为郑忧矣。﹂文公弗听。\\
\hspace*{2em}三十七年春。晋公子重耳反国。立。是为文公。秋。郑入滑。滑听命。已而反与卫。於是郑伐滑。周襄王使伯馃请滑。郑文公怨惠王之亡在栎。而文公父厉公入之。而惠王不赐厉公爵禄。又怨襄王之与卫滑。故不听襄王请而囚伯馃。王怒。与翟人伐郑。弗克。冬。翟攻伐襄王。襄王出奔郑。郑文公居王于氾。三十八年。晋文公入襄王成周。\\
\hspace*{2em}四十一年。助楚击晋。自晋文公之过无礼。故背晋助楚。四十三年。晋文公与秦穆公共围郑。讨其助楚攻晋者。及文公过时之无礼也。初。郑文公有三夫人。宠子五人。皆以罪蚤死。公怒。溉逐群公子。子兰奔晋。从晋文公围郑。时兰事晋文公甚谨。爱幸之。乃私於晋。以求入郑为太子。晋於是欲得叔詹为僇。郑文公恐。不敢谓叔詹言。詹闻。言於郑君曰。﹁臣谓君。君不听臣。晋卒为患。然晋所以围郑。以詹。詹死而赦郑国。詹之原也。﹂乃自杀。郑人以詹尸与晋。晋文公曰。﹁必欲一见郑君。辱之而去。﹂郑人患之。乃使人私於秦曰。﹁破郑益晋。非秦之利也。﹂秦兵罢。晋文公欲入兰为太子。以告郑。郑大夫石癸曰。﹁吾闻姞姓乃后稷之元妃。其後当有兴者。子兰母。其後也。且夫人子尽已死。馀庶子无如兰贤。今围急。晋以为请。利孰大焉。﹂遂许晋。与盟。而卒立子兰为太子。晋兵乃罢去。\\
\hspace*{2em}四十五年。文公卒。子兰立。是为缪公。\\
\hspace*{2em}缪公元年春。秦缪公使三将将兵欲袭郑。至滑。逢郑贾人弦高诈以十二牛劳军。故秦兵不至而还。晋败之於崤。初。往年郑文公之卒也。郑司城缯贺以郑情卖之。秦兵故来。三年。郑发兵从晋伐秦。败秦兵於汪。\\
\hspace*{2em}往年楚太子商臣弑其父成王代立。二十一年。与宋华元伐郑。华元杀羊食士。不与其御羊斟。怒以驰郑。郑囚华元。宋赎华元。元亦亡去。晋使赵穿以兵伐郑。\\
\hspace*{2em}二十二年。郑缪公卒。子夷立。是为灵公。\\
\hspace*{2em}灵公元年春。楚献鼋於灵公。子家。子公将朝灵公。子公之食指动。谓子家曰。﹁佗日指动。必食异物。﹂及入。见灵公进鼋羹。子公笑曰。﹁果然。﹂灵公问其笑故。具告灵公。灵公召之。独弗予羹。子公怒。染其指。尝之而出。公怒。欲杀子公。子公与子家谋先。夏。弑灵公。郑人欲立灵公弟去疾。去疾让曰。﹁必以贤。则去疾不肖。必以顺。则公子坚长。﹂坚者。灵公庶弟。去疾之兄也。於是乃立子坚。是为襄公。\\
\hspace*{2em}襄公立。将尽去缪氏。缪氏者。杀灵公。子公之族家也。去疾曰。﹁必去缪氏。我将去之。﹂乃止。皆以为大夫。\\
\hspace*{2em}襄公元年。楚怒郑受宋赂纵华元。伐郑。郑背楚。与晋亲。五年。楚复伐郑。晋来救之。六年。子家卒。国人复逐其族。以其弑灵公也。\\
\hspace*{2em}七年。郑与晋盟鄢陵。八年。楚庄王以郑与晋盟。来伐。围郑三月。郑以城降楚。楚王入自皇门。郑襄公肉袒掔羊以迎。曰。﹁孤不能事边邑。使君王怀怒以及弊邑。孤之罪也。敢不惟命是听。君王迁之江南。及以赐诸侯。亦惟命是听。若君王不忘厉。宣王。桓。武公。哀不忍绝其社稷。锡不毛之地。使复得改事君王。孤之原也。然非所敢望也。敢布腹心。惟命是听。﹂庄王为卻三十里而後舍。楚群臣曰。﹁自郢至此。士大夫亦久劳矣。今得国舍之。何如。﹂庄王曰。﹁所为伐。伐不服也。今已服。尚何求乎。﹂卒去。晋闻楚之伐郑。发兵救郑。其来持两端。故迟。比至河。楚兵已去。晋将率或欲渡。或欲还。卒渡河。庄王闻。还击晋。郑反助楚。大破晋军於河上。十年。晋来伐郑。以其反晋而亲楚也。\\
\hspace*{2em}十一年。楚庄王伐宋。宋告急于晋。晋景公欲发兵救宋。伯宗谏晋君曰。﹁天方开楚。未可伐也。﹂乃求壮士得霍人解扬。字子虎。诓楚。令宋毋降。过郑。郑与楚亲。乃执解扬而献楚。楚王厚赐与约。使反其言。令宋趣降。三要乃许。於是楚登解扬楼车。令呼宋。遂负楚约而致其晋君命曰。﹁晋方悉国兵以救宋。宋虽急。慎毋降楚。晋兵今至矣。﹂楚庄王大怒。将杀之。解扬曰。﹁君能制命为义。臣能承命为信。受吾君命以出。有死无陨。﹂庄王曰。﹁若之许我。已而背之。其信安在。﹂解扬曰。﹁所以许王。欲以成吾君命也。﹂将死。顾谓楚军曰。﹁为人臣无忘尽忠得死者。﹂楚王诸弟皆谏王赦之。於是赦解扬使归。晋爵之为上卿。\\
\hspace*{2em}十八年。襄公卒。子悼公晞立。\\
\hspace*{2em}悼公元年。鄦公恶郑於楚。悼公使弟睔於楚自讼。讼不直。楚囚睔。於是郑悼公来与晋平。遂亲。睔私於楚子反。子反言归睔於郑。\\
\hspace*{2em}二年。楚伐郑。晋兵来救。是岁。悼公卒。立其弟睔。是为成公。\\
\hspace*{2em}成公三年。楚共王曰﹁郑成公孤有德焉﹂。使人来与盟。成公私与盟。秋。成公朝晋。晋曰﹁郑私平於楚﹂。执之。使栾书伐郑。四年春。郑患晋围。公子如乃立成公庶兄繻为君。其四月。晋闻郑立君。乃归成公。郑人闻成公归。亦杀君繻。迎成公。晋兵去。\\
\hspace*{2em}十年。背晋盟。盟於楚。晋厉公怒。发兵伐郑。楚共王救郑。晋楚战鄢陵。楚兵败。晋射伤楚共王目。俱罢而去。十三年。晋悼公伐郑。兵於洧上。郑城守。晋亦去。\\
\hspace*{2em}十四年。成公卒。子恽立。是为釐公。\\
\hspace*{2em}釐公五年。郑相子驷朝釐公。釐公不礼。子驷怒。使厨人药杀釐公。赴诸侯曰﹁釐公暴病卒﹂。立釐公子嘉。嘉时年五岁。是为简公。\\
\hspace*{2em}简公元年。诸公子谋欲诛相子驷。子驷觉之。反尽诛诸公子。二年。晋伐郑。郑与盟。晋去。冬。又与楚盟。子驷畏诛。故两亲晋。楚。三年。相子驷欲自立为君。公子子孔使尉止杀相子驷而代之。子孔又欲自立。子产曰。﹁子驷为不可。诛之。今又效之。是乱无时息也。﹂於是子孔从之而相郑简公。\\
\hspace*{2em}四年。晋怒郑与楚盟。伐郑。郑与盟。楚共王救郑。败晋兵。简公欲与晋平。楚又囚郑使者。\\
\hspace*{2em}十二年。简公怒相子孔专国权。诛之。而以子产为卿。十九年。简公如晋请卫君还。而封子产以六邑。子产让。受其三邑。二十二年。吴使延陵季子於郑。见子产如旧交。谓子产曰。﹁郑之执政者侈。难将至。政将及子。子为政。必以礼。不然。郑将败。﹂子产厚遇季子。二十三年。诸公子争宠相杀。又欲杀子产。公子或谏曰。﹁子产仁人。郑所以存者子产也。勿杀。﹂乃止。\\
\hspace*{2em}二十五年。郑使子产於晋。问平公疾。平公曰。﹁卜而曰实沈。台骀为祟。史官莫知。敢问。﹂对曰。﹁高辛氏有二子。长曰阏伯。季曰实沈。居旷林。不相能也。日操干戈以相征伐。后帝弗臧。迁阏伯于商丘。主辰。商人是因。故辰为商星。迁实沈于大夏。主参。唐人是因。服事夏。商。其季世曰唐叔虞。﹂当武王邑姜方娠大叔。梦帝谓己。﹃余命而子曰虞。乃与之唐。属之参而蕃育其子孙。﹄及生有文在其掌曰﹃虞﹄。遂以命之。及成王灭唐而国大叔焉。故参为晋星。﹂由是观之。则实沈。参神也。昔金天氏有裔子曰昧。为玄冥师。生允格。台骀。台骀能业其官。宣汾。洮。障大泽。以处太原。帝用嘉之。国之汾川。沈。姒。蓐。黄实守其祀。今晋主汾川而灭之。由是观之。则台骀。汾。洮神也。然是二者不害君身。山川之神。则水旱之菑禜之。日月星辰之神。则雪霜风雨不时禜之。若君疾。饮食哀乐女色所生也。﹂平公及叔乡曰。﹁善。博物君子也。﹂厚为之礼於子产。\\
\hspace*{2em}二十七年夏。郑简公朝晋。冬。畏楚灵王之彊。又朝楚。子产从。二十八年。郑君病。使子产会诸侯。与楚灵王盟於申。诛齐庆封。\\
\hspace*{2em}三十六年。简公卒。子定公宁立。秋。定公朝晋昭公。\\
\hspace*{2em}定公元年。楚公子弃疾弑其君灵王而自立。为平王。欲行德诸侯。归灵王所侵郑地于郑。\\
\hspace*{2em}四年。晋昭公卒。其六卿彊。公室卑。子产谓韩宣子曰。﹁为政必以德。毋忘所以立。﹂\\
\hspace*{2em}六年。郑火。公欲禳之。子产曰。﹁不如修德。﹂\\
\hspace*{2em}八年。楚太子建来奔。十年。太子建与晋谋袭郑。郑杀建。建子胜奔吴。\\
\hspace*{2em}十一年。定公如晋。晋与郑谋。诛周乱臣。入敬王于周。\\
\hspace*{2em}十三年。定公卒。子献公虿立。献公十三年卒。子声公胜立。当是时。晋六卿彊。侵夺郑。郑遂弱。\\
\hspace*{2em}声公五年。郑相子产卒。郑人皆哭泣。悲之如亡亲戚。子产者。郑成公少子也。为人仁爱人。事君忠厚。孔子尝过郑。与子产如兄弟云。及闻子产死。孔子为泣曰。﹁古之遗爱也。﹂\\
\hspace*{2em}八年。晋范。中行氏反晋。告急於郑。郑救之。晋伐郑。败郑军於铁。\\
\hspace*{2em}十四年。宋景公灭曹。二十年。齐田常弑其君简公。而常相於齐。二十二年。楚惠王灭陈。孔子卒。\\
\hspace*{2em}三十六年。晋知伯伐郑。取九邑。\\
\hspace*{2em}三十七年。声公卒。子哀公易立。哀公八年。郑人弑哀公而立声公弟丑。是为共公。共公三年。三晋灭知伯。三十一年。共公卒。子幽公已立。幽公元年。韩武子伐郑。杀幽公。郑人立幽公弟骀。是为繻公。\\
\hspace*{2em}繻公十五年。韩景侯伐郑。取雍丘。郑城京。\\
\hspace*{2em}十六年。郑伐韩。败韩兵於负黍。二十年。韩。赵。魏列为诸侯。二十三年。郑围韩之阳翟。\\
\hspace*{2em}二十五年。郑君杀其相子阳。二十七。子阳之党共弑繻公骀而立幽公弟乙为君。是为郑君。\\
\hspace*{2em}郑君乙立二年。郑负黍反。复归韩。十一年。韩伐郑。取阳城。\\
\hspace*{2em}二十一年。韩哀侯灭郑。并其国。\\
\hspace*{2em}太史公曰。语有之。﹁以权利合者。权利尽而交疏﹂。甫瑕是也。甫瑕虽以劫杀郑子内厉公。厉公终背而杀之。此与晋之里克何异。守节如荀息。身死而不能存奚齐。变所从来。亦多故矣。\\

% ===== 索隐述赞（司马贞《史记索隐》）· 统一样式 =====
% 题头「索隐述赞曰」五字并入反白框；正文每四句一列，定宽盒 \hfil 居列中；
% \fontsize{24}{41} 列距对齐全书界行（121+41k）；\clearpage 使述赞独占一页。
\clearpage
\shuzanhead{索隐述赞曰}\linebreak
{\shuzanbody\fontsize{24}{41}\selectfont
\hbox to \hsize{\hfil 厉王之子。得封於郑。代职司徒。缁衣在咏。\hfil}\linebreak
\hbox to \hsize{\hfil 虢。郐献邑。祭祝专命。庄既犯王。\hfil}\linebreak
\hbox to \hsize{\hfil 厉亦奔命。居栎克入。梦兰毓庆。伯服生囚。\hfil}\linebreak
\hbox to \hsize{\hfil 叔瞻尸聘。釐。简之後。公室不竞。\hfil}\linebreak
\hbox to \hsize{\hfil 负黍虽还。韩哀日盛。\hfil}%
\par}
\end{document}
